Write $\neg \eventuallyZero$ for the \emph{eventually one} sequences,
the complements of those in $\eventuallyZero$,
and $\eventuallyConstant := \eventuallyZero \cup \neg\eventuallyZero$ for the eventually
constant ones.

\begin{defi}\label{def.twosComplementSequence}
 Let $\twosComplement : \Z \to \bitSequences$ be
 \begin{equation*}
  \twosComplement(x) :=
  \begin{cases}
   \binaryRepresentation(x), & x \geq 0, \\
   \neg\, \binaryRepresentation(-1 - x), & x < 0 .
  \end{cases}
 \end{equation*}
\end{defi}

\begin{prop}\label{prop.integersAsSequences}
 $\twosComplement$ is injective, with image $\eventuallyConstant$, and for every $n$ and every
 $x \in \Z$,
 \begin{equation}\label{eq.integersAsSequences}
  \truncation{n}\big(\twosComplement(x)\big) = \digitsOf{n}\big(\reduction{n}(x)\big) .
 \end{equation}
\end{prop}
\begin{proof}
 For $x \geq 0$, $\twosComplement(x) = \binaryRepresentation(x) \in \eventuallyZero$;
 for $x < 0$, $-1 - x \geq 0$, so $\binaryRepresentation(-1-x)$ is eventually zero and
 $\twosComplement(x)$ eventually one.
 No sequence is both, so the two cases have disjoint images;
 each is injective by Proposition \ref{prop.binaryRepresentation},
 and together their images are $\eventuallyZero \cup \neg\eventuallyZero$.
 For $x \geq 0$, \eqref{eq.integersAsSequences} is Proposition \ref{prop.truncationCommutes}.
 For $x < 0$ put $m = -1 - x$.
 Truncating $\neg\,\binaryRepresentation(m)$ complements the truncation of
 $\binaryRepresentation(m)$,
 whose value is $\reduction{n}(m)$ by Proposition \ref{prop.truncationCommutes};
 by Proposition \ref{prop.twosComplement} the complement has the value
 $(2^{n+1} - 1) - \reduction{n}(m) \equiv -1 - m = x$ modulo $2^{n+1}$,
 and a value of $n+1$ bits congruent to $x$ is $\reduction{n}(x)$.
\end{proof}

$\twosComplement$ is not onto $\bitSequences$:
the alternating sequence $0, 1, 0, 1, \dots$ is not eventually constant.

The set $\eventuallyConstant$ is closed under $\neg$, $\wedge$, $\vee$ and $\oplus$,
since beyond the position where both arguments have become constant each of them returns a
constant bit.
So the four operators carry over to the integers,
and on the integers they are the bitwise operators of Python:
\codehl{\textasciitilde x} is $\twosComplement^{-1}(\neg\,\twosComplement(x))$,
\codehl{x \& y} is $\twosComplement^{-1}(\twosComplement(x) \wedge \twosComplement(y))$,
and so on.
We write $x \bitAnd y$ and $x \bitOr y$ for the two that Chapter \ref{chap.microstructure}
could misread, since there $\wedge$ and $\vee$ between two integers are their minimum and
maximum;
$\neg x$ and $x \oplus y$ carry no such meaning, and keep their symbols.
By Definition \ref{def.twosComplementSequence}, $\neg x = -x - 1$ on the whole of $\Z$.
Python computes each operator in one pass over the digits it stores:
nothing infinite is built,
and \codehl{bits \&= \textasciitilde bit} of Listing \ref{listing.tickArrayWrite},
which meets the leading ones of $\neg\,$\codehl{bit} with the leading zeros of \codehl{bits},
returns a non-negative integer.

Bit $k$ of an integer $x$ is \codehl{x >> k \& 1}, negative $x$ included,
because \codehl{x >> k} is $\lfloor x / 2^k \rfloor$:
this is $\twosComplement$, computed.
Listing \ref{listing.bitsOfIntegers} checks both facts on a thousand integers of either sign.

\begin{lstlisting}[caption={$\neg x = -x-1$, and the bits of an integer of either sign.}, label=listing.bitsOfIntegers]
import random

for _ in range(1000):
    x = random.randint(-(2**64), 2**64)
    assert ~x == -x - 1
    assert all((x >> k & 1) == (x // 2**k) % 2 for k in range(70))
\end{lstlisting}

For $b = 360$ the shifts reach a fixed point on either sign, as Table \ref{tab.shiftsOf360}
shows: $0$ for $b$, whose bits are all zero beyond position $8$, and $-1$ for $-b$, whose bits
are all one beyond position $8$.
So the infinite tails are generated rather than stored,
\begin{equation*}
 \twosComplement(360) = \dots 0000\,101101000,
 \qquad
 \twosComplement(-360) = \dots 1111\,010011000 ,
\end{equation*}
written with position $0$ on the right.
Position $3$ is the only one where both have a $1$,
and the next section shows that this is always so.

\begin{table}[h!]
 \centering
 \begin{tabular}{@{}rrcrc@{}}
  \toprule
  $k$ & $\lfloor 360/2^k \rfloor$ & bit & $\lfloor -360/2^k \rfloor$ & bit \\
  \midrule
  $0$ & $360$ & $0$ & $-360$ & $0$ \\
  $1$ & $180$ & $0$ & $-180$ & $0$ \\
  $2$ & $90$ & $0$ & $-90$ & $0$ \\
  $3$ & $45$ & $1$ & $-45$ & $1$ \\
  $4$ & $22$ & $0$ & $-23$ & $1$ \\
  $5$ & $11$ & $1$ & $-12$ & $0$ \\
  $6$ & $5$ & $1$ & $-6$ & $0$ \\
  $7$ & $2$ & $0$ & $-3$ & $1$ \\
  $8$ & $1$ & $1$ & $-2$ & $0$ \\
  $9$ & $0$ & $0$ & $-1$ & $1$ \\
  $10$ & $0$ & $0$ & $-1$ & $1$ \\
  \bottomrule
 \end{tabular}
 \caption{The shifts of $360$ and of $-360$, and the bit each leaves at position $0$.}
 \label{tab.shiftsOf360}
\end{table}
